\chapter*{Foreword}
\addcontentsline{toc}{chapter}{Foreword}

Model Predictive Control (MPC) is an optimisation-based control methodology. By explicitly exploiting a dynamical model of the system, MPC predicts future behaviour and determines control actions by solving a constrained optimisation problem in real time. This distinguishes MPC from classical feedback controllers such as PID or lead--lag compensators, which rely on fixed control laws designed offline.

Initially developed in the process industries~\cite{Qin2003}, MPC has since been adopted in robotics (trajectory tracking under actuator limits), autonomous vehicles (path planning with obstacle avoidance), and aerospace (powered descent guidance under thrust constraints). Its ability to handle multivariable dynamics and hard constraints within a single optimisation framework makes it applicable wherever classical single-loop controllers reach their limits.

\paragraph{Aim and scope.}
This book has been written with a strongly practical emphasis. The primary objective is to introduce MPC in a way that allows readers to apply the techniques they learn directly to realistic engineering problems, rather than treating MPC as a purely theoretical topic. Wherever possible, abstract concepts are motivated using physical intuition, illustrative examples, and implementation-oriented discussions. Mathematical depth is deliberately balanced against clarity and applicability throughout.

The first eight chapters form the core of the book. They take the reader from the mathematical foundations---optimisation, state-space modelling, controllability, observability---through the LQR and Kalman filter, to the formulation and implementation of constrained linear MPC with output feedback. Each chapter includes worked examples and MATLAB code so that readers can reproduce every result.

Beyond the core, several supplementary chapters extend MPC into more advanced territory: robust and tube MPC for systems with bounded disturbances, nonlinear MPC via direct transcription and sequential convexification, and learning-based MPC methods including GP-MPC, Koopman MPC, and data-driven predictive control (DeePC). A Simulink implementation chapter demonstrates closed-loop MPC in a building energy system. Two self-study appendices review the required linear algebra and the probability theory used in the Kalman filter chapters.

The book is structured so that it can be followed either as a companion to a taught course or as a self-contained resource for independent study. A reader who works through the chapters sequentially should acquire a working knowledge of MPC sufficient to model systems, formulate predictive control problems, and implement controllers using modern computational tools.

\paragraph{The $^{(*)}$ symbol.} Chapters and sections marked with $^{(*)}$ in their title contain supplementary material that goes beyond the core exposition. This includes the formal stability theory section in the MPC chapter (recursive feasibility, terminal sets, and the Lyapunov stability proof), the Simulink implementation chapter (Chapter~\ref{chap:simulink_mpc}), the Robust and Tube MPC chapter (Chapter~\ref{chap:tube_mpc}), the Nonlinear MPC chapter (Chapter~\ref{chap:nmpc}), the MPC with Learning chapter (Chapter~\ref{chap:learning_mpc}), and the conclusion chapter (Chapter~\ref{chap:conclusion}). The core material is contained in the unmarked chapters and sections.

\paragraph{Software prerequisites.}
The MATLAB code throughout this book requires:
\begin{itemize}[noitemsep]
  \item MATLAB R2019b or later (for \texttt{yline}, \texttt{tiledlayout}).
  \item \emph{Control System Toolbox} --- \texttt{ctrb}, \texttt{obsv}, \texttt{dlqr}, \texttt{dare}, \texttt{c2d}, \texttt{ss}.
  \item \emph{Optimization Toolbox} --- \texttt{quadprog}, \texttt{linprog}.
  \item \emph{YALMIP} (free, \url{https://yalmip.github.io}) --- used in every MPC chapter.
  \item \emph{MPT3} (free, \url{https://www.mpt3.org}) --- polytopic sets in Chapters~\ref{chap:linear_mpc} and~\ref{chap:tube_mpc}.
  \item \emph{CasADi} (free, \url{https://web.casadi.org}) --- nonlinear MPC in Chapter~\ref{chap:nmpc}.
\end{itemize}
Optional toolboxes: Statistics and Machine Learning Toolbox (\texttt{normpdf}, \texttt{normcdf}, \texttt{norminv}, \texttt{mvnrnd}), Deep Learning Toolbox (\texttt{feedforwardnet}), and Simulink.

\paragraph{Prerequisites.}
A basic background in SISO control systems, state-space representations, linear algebra, and numerical computation is assumed. Familiarity with MATLAB is strongly recommended, as examples and implementations throughout the book rely on MATLAB-based workflows. With this background, the material should be accessible to graduate-level readers in control, instrumentation, and related engineering disciplines. Readers who need to refresh specific mathematical topics are encouraged to consult the appendices as needed.

\vspace{1.5cm}
\begin{flushright}
\textit{\.Ibrahim K\"u\c{c}\"ukdemiral}\\
Glasgow, Scotland, UK\\
2026
\end{flushright}
